% latex/supplement/main.tex
%
% Master document for the A3E-DNS supplement.
%
% The supplement is intended to be typeset separately from the main
% thesis and included in the thesis's appendix as a separate PDF. It
% contains the experiment descriptions, the notation table, the figure
% captions, and the generated result tables. The result tables are
% \input{} rather than inlined so that re-running the generators
% updates the supplement without a manual edit.

\documentclass[11pt,a4paper]{article}

% ----- Packages ---------------------------------------------------- %
\usepackage[T1]{fontenc}
\usepackage[utf8]{inputenc}
\usepackage{lmodern}
\usepackage{microtype}
\usepackage[margin=2.5cm]{geometry}
\usepackage{amsmath,amssymb,amsthm}
\usepackage{booktabs}
\usepackage{graphicx}
\usepackage{siunitx}
\usepackage{hyperref}
\usepackage{xcolor}
\usepackage{longtable}
\usepackage{enumitem}
\usepackage{caption}
\usepackage{subcaption}
\usepackage{float}
\usepackage{titlesec}

% ----- Hyperref setup ---------------------------------------------- %
\hypersetup{
    colorlinks=true,
    linkcolor=blue,
    citecolor=blue,
    urlcolor=blue,
    pdftitle={A3E-DNS Supplement},
    pdfauthor={A3E-DNS Extended Suite},
}

% ----- SIunitx setup ----------------------------------------------- %
\sisetup{
    detect-all,
    group-separator={,},
    group-minimum-digits=4,
    round-mode=figures,
    round-precision=3,
}

% ----- Theorem-like environments ----------------------------------- %
\newtheorem{proposition}{Proposition}
\newtheorem{definition}{Definition}
\newtheorem{lemma}{Lemma}

% ----- Caption setup ----------------------------------------------- %
\captionsetup{
    font=small,
    labelfont=bf,
    justification=raggedright,
    singlelinecheck=false,
}

% ----- Title ------------------------------------------------------- %
\title{
    \textbf{Supplementary Material} \\
    \large Next-Generation Secure Encrypted DNS Architectures: \\
    \large Four Complementary Frameworks for Adaptive, \\
    \large Privacy-Preserving, and Quantum-Resilient Defense
}
\author{A3E-DNS Extended Suite}
\date{November 30, 2026}

% ----- Shared macros ----------------------------------------------- %
% latex/supplement/macros.tex
%
% Shared macros for the A3E-DNS supplement. Every macro is defined
% once here and re-used throughout the supplement's chapters so that a
% change to, for example, the DP epsilon symbol propagates everywhere.

% ----- The four frameworks ------------------------------------------ %
\newcommand{\mirage}{\textsc{Mirage-DoH}}
\newcommand{\pafa}{\textsc{Pafa-DoH}}
\newcommand{\eThree}{E3-DoH}
\newcommand{\aThreeE}{A\textsuperscript{3}E-DNS}

% ----- Protocols ---------------------------------------------------- %
\newcommand{\DoH}{\textsc{DoH}}
\newcommand{\DoT}{\textsc{DoT}}
\newcommand{\DoQ}{\textsc{DoQ}}

% ----- Cryptographic symbols ---------------------------------------- %
\newcommand{\Qring}{\mathbb{Z}_Q}
\newcommand{\Ndim}{N}
\newcommand{\ntt}{\mathsf{NTT}}
\newcommand{\intt}{\mathsf{NTT}^{-1}}
\newcommand{\sigmas}{{\sigma_{\mathrm{s}}}}
\newcommand{\eps}{\varepsilon}

% ----- AEE state ---------------------------------------------------- %
\newcommand{\se}{\mathfrak{se}(3)}
\newcommand{\SE}{\mathrm{SE}(3)}
\newcommand{\xivec}{\xi}
\newcommand{\rhovec}{\rho}
\newcommand{\omegavec}{\omega}
\newcommand{\psivec}{\psi}
\newcommand{\volatility}{v}

% ----- Byzantine / FL ----------------------------------------------- %
\newcommand{\nclients}{n}
\newcommand{\nbyzantine}{f}
\newcommand{\agg}{\mathsf{Agg}}
\newcommand{\med}{\mathsf{med}}
\newcommand{\trim}{\mathsf{trim}}

% ----- Metrics ------------------------------------------------------ %
\newcommand{\acc}{\mathsf{Acc}}
\newcommand{\fpr}{\mathsf{FPR}}
\newcommand{\pnn}[1]{P_{#1}}

% ----- Thesis reference markers ------------------------------------- %
\newcommand{\thesisref}[1]{\textcolor{gray}{\texttt{[#1]}}}
\newcommand{\TBD}[1]{\textcolor{red}{TODO: #1}}

% ----- Convenience -------------------------------------------------- %
\newcommand{\eg}{e.g.}
\newcommand{\ie}{i.e.}
\newcommand{\etal}{\textit{et al.}}


\begin{document}

\maketitle

\begin{abstract}
\noindent
This supplement documents the extended experimentation suite that
accompanies Chapter~6 of the thesis \emph{Next-Generation Secure
Encrypted DNS Architectures}. It contains the descriptions of the
ten experiments, the notation table, the figure captions, and the
generated result tables. Every number reported here is produced by
the code in \texttt{extended\_experimentation/}; the supplement
itself does not recompute anything. A reviewer who runs the pipeline
against the same seeds should reproduce the tables to the precision
shown.
\end{abstract}

\tableofcontents
\newpage

% ----- Sections ---------------------------------------------------- %
% latex/supplement/notation.tex
%
% Symbol table for the A3E-DNS supplement. The table is generated by
% hand rather than from a script because the descriptions are
% load-bearing: a reviewer consulting the table needs to see what
% each symbol is used for, not only its spelling.

\section*{Notation}

The following table lists the symbols used throughout this supplement.
Symbols are grouped by the subsystem they belong to; a symbol that
appears in more than one subsystem is defined in the earlier group and
cross-referenced in the later one.

\subsection*{Protocols and frameworks}

\begin{center}
\begin{tabular}{ll}
\toprule
Symbol & Meaning \\
\midrule
$\DoH$, $\DoT$, $\DoQ$ & DNS-over-HTTPS, DNS-over-TLS, DNS-over-QUIC \\
$\mirage$              & First framework: metamorphic defence \\
$\pafa$                & Second framework: federated privacy \\
$\eThree$              & Third framework: multi-protocol expansion \\
$\aThreeE$             & Fourth framework: adaptive evolutionary encryption \\
\bottomrule
\end{tabular}
\end{center}

\subsection*{Ring arithmetic (QHE)}

\begin{center}
\begin{tabular}{ll}
\toprule
Symbol & Meaning \\
\midrule
$Q$              & The ML-DSA prime, $Q = 8\,380\,417$ \\
$\Ndim$          & Ring dimension, $\Ndim = 256$ \\
$\Qring$         & $\mathbb{Z}_Q$, the finite field of integers modulo $Q$ \\
$\Qring[X]/(X^{\Ndim}+1)$ & The negacyclic ring the NTT operates over \\
$\psi$           & A primitive $2\Ndim$-th root of unity modulo $Q$ \\
$\omega$         & $\psi^2$, a primitive $\Ndim$-th root of unity \\
$\ntt$           & The negacyclic NTT operator \\
$\intt$          & Its inverse \\
$\sigmas$        & Sampler standard deviation, $\sigmas = 3.2$ \\
$\eps$           & Differential-privacy budget, $\eps = 0.68$ \\
\bottomrule
\end{tabular}
\end{center}

\subsection*{Adaptive evolutionary encryption (AEE)}

\begin{center}
\begin{tabular}{ll}
\toprule
Symbol & Meaning \\
\midrule
$\se$            & The Lie algebra of $\SE$, spanned by $(\rhovec,\omegavec)$ \\
$\SE$            & The special Euclidean group of rigid motions \\
$\xivec$         & AEE state, $\xivec = (\rhovec,\omegavec) \in \mathbb{R}^6$ \\
$\volatility$    & Volatility scalar, floored at $\volatility_{\min} = 0.05$ \\
$\alpha$         & Diffusion scale, $\alpha = 0.12$ \\
$\beta$          & Drift scale, $\beta = -0.02$ \\
\bottomrule
\end{tabular}
\end{center}

\subsection*{Byzantine / federated learning}

\begin{center}
\begin{tabular}{ll}
\toprule
Symbol & Meaning \\
\midrule
$\nclients$      & Total number of clients in a round \\
$\nbyzantine$    & Number of Byzantine clients \\
$\agg$           & Aggregation operator \\
$\med$           & Coordinate-wise median aggregator \\
$\trim$          & Coordinate-wise trimmed-mean aggregator \\
$\acc$           & Detection accuracy \\
$\fpr$           & False-positive rate, as a fraction \\
$\pnn{q}$        & The $q$-th percentile of a latency distribution \\
\bottomrule
\end{tabular}
\end{center}

\subsection*{Thesis reference values}

The following values are the headline results of the A3E-DNS
evaluation. Every table in this supplement that reports one of these
metrics can be cross-checked against the value below; a reviewer
who finds a table that does not match should treat that as a bug in
this supplement, not in the thesis.

\begin{center}
\begin{tabular}{lr}
\toprule
Metric & Thesis value \\
\midrule
Detection accuracy              & $99.94\%$ \\
False-positive rate             & $0.17\%$ \\
P99 latency                     & $1.42\,\mathrm{ms}$ \\
Peak throughput                 & $22.5\,\mathrm{M}\,\mathrm{QPS}$ \\
Adversarial robustness          & $96.5\%$ \\
Differential-privacy budget     & $\eps = 0.68$ \\
\bottomrule
\end{tabular}
\end{center}

\newpage

% latex/supplement/experiments.tex
%
% Descriptions of the ten experiments. For each experiment the
% supplement records the purpose, the inputs, the artefact paths, and
% the thesis table the results feed. A reviewer running the pipeline
% should be able to identify every output directory mentioned here
% by name.

\section{Experiment descriptions}

This section documents the ten experiments the A3E-DNS evaluation
reports. Each experiment is a standalone Python script under
\texttt{experiments\_v2/}, and each is invoked by the orchestrator in
\texttt{extended\_experimentation/experiments/e1\_e10\_orchestrator.py}.
The descriptions below follow the order the orchestrator runs them in.

\subsection{E1 --- Attack scenario catalogue}

\textbf{Purpose.} Produce the ten-scenario attack catalogue with
three severity variants per scenario and three protocols per scenario,
plus a benign baseline.

\textbf{Inputs.} None; the catalogue is fully synthetic, parameterised
by the seed.

\textbf{Outputs.}
\begin{itemize}
\item \texttt{runs/experiments/E1/scenario\_sweep.json} — per-scenario
      accuracy, FPR, precision, recall.
\item \texttt{runs/experiments/E1/predictions.jsonl} — per-window
      predictions.
\item \texttt{runs/experiments/E1/manifest.json} — the run manifest.
\end{itemize}

\textbf{Feeds.} Table 30 (scenario breakdown), Table 38 (adversarial
robustness), Table 39 (statistical significance).

\subsection{E2 --- Byzantine robustness sweep}

\textbf{Purpose.} Sweep the Byzantine fraction from $5\%$ to $40\%$
across the eight attack strategies and the four schedules.

\textbf{Inputs.} The E1 predictions (as the honest baseline).

\textbf{Outputs.}
\begin{itemize}
\item \texttt{runs/experiments/E2/byzantine\_sweep.json} — per-round
      attack ratios and aggregate accuracies.
\item \texttt{runs/experiments/E2/manifest.json}.
\end{itemize}

\textbf{Feeds.} Table 31 (Byzantine robustness). The
\texttt{attack\_to\_reference\_ratio} at each fraction is what the
robustness figure plots.

\subsection{E3 --- Differential-privacy sweep}

\textbf{Purpose.} Sweep the target $\eps$ from $0.1$ to $1.0$ and
measure the utility loss at each value.

\textbf{Inputs.} The E1 predictions.

\textbf{Outputs.}
\begin{itemize}
\item \texttt{runs/experiments/E3/dp\_sweep.json} — per-$\eps$
      accuracy and FPR.
\item \texttt{runs/experiments/E3/manifest.json}.
\end{itemize}

\textbf{Feeds.} Table 32 (DP epsilon sweep). The reported
$\eps = 0.68$ is the point at which the utility loss first crosses
$0.5\%$, which the sweep identifies from the data rather than from a
pre-set threshold.

\subsection{E4a--E4d --- QHE primitives}

\textbf{Purpose.} Four separate benchmarks of the homomorphic
encryption primitives:

\begin{itemize}
\item E4a: negacyclic NTT vs.\ naive polynomial multiplication;
\item E4b: CDT Gaussian sampler correctness and timing;
\item E4c: selective homomorphic operations over the informative
      feature subset;
\item E4d: packed ciphertext size and compression ratio.
\end{itemize}

\textbf{Inputs.} All four are self-contained; none depends on the
attack catalogue.

\textbf{Outputs.} Four JSON manifests under
\texttt{runs/experiments/E4a/} through \texttt{E4d/}.

\textbf{Feeds.} Tables 33 and 34.

\subsection{E5 --- End-to-end P99 latency}

\textbf{Purpose.} Measure the P99 latency of the full pipeline under
combined attack load.

\textbf{Inputs.} The E2 and E3 artefacts, replayed through the
pipeline.

\textbf{Outputs.}
\begin{itemize}
\item \texttt{runs/experiments/E5/latency.json} — per-protocol P50,
      P95, P99, P99.9.
\item \texttt{runs/experiments/E5/manifest.json}.
\end{itemize}

\textbf{Feeds.} Table 35 (P99 latency). The reported
$1.42\,\mathrm{ms}$ is the P99 over the DoH, DoT and DoQ streams
pooled; the per-protocol values are in the same table as a
footnote.

\subsection{E6 --- Peak throughput}

\textbf{Purpose.} Measure peak QPS at concurrency levels from 1 to
128.

\textbf{Inputs.} The E5 pipeline configuration.

\textbf{Outputs.}
\begin{itemize}
\item \texttt{runs/experiments/E6/throughput.json} — per-concurrency
      QPS.
\item \texttt{runs/experiments/E6/manifest.json}.
\end{itemize}

\textbf{Feeds.} Table 35 (peak throughput). The reported
$22.5\,\mathrm{M}\,\mathrm{QPS}$ is the maximum over the sweep,
measured at concurrency $= 32$.

\subsection{E7 --- Component ablation}

\textbf{Purpose.} Ablate each of AEE, QHE, FL and the attention block
in turn, and measure the resulting accuracy, FPR and P99.

\textbf{Inputs.} Pre-computed checkpoints for the four ablated
variants, produced by re-running the training pipeline with the
corresponding component disabled.

\textbf{Outputs.}
\begin{itemize}
\item \texttt{runs/experiments/E7/ablation.json}.
\item \texttt{runs/experiments/E7/manifest.json}.
\end{itemize}

\textbf{Feeds.} Table 36. The ablation is \emph{partial}: a full
factorial design was not performed, as Section~6.7.3 of the thesis
states.

\subsection{E8 --- Multi-protocol accuracy}

\textbf{Purpose.} Confirm the detector generalises across DoH, DoT
and DoQ.

\textbf{Inputs.} The E1 catalogue, stratified by protocol.

\textbf{Outputs.}
\begin{itemize}
\item \texttt{runs/experiments/E8/multi\_protocol.json}.
\item \texttt{runs/experiments/E8/manifest.json}.
\end{itemize}

\textbf{Feeds.} Table 37.

\subsection{E9 --- Adversarial robustness}

\textbf{Purpose.} Measure the accuracy drop under the adaptive-blend
variant of each scenario, relative to the baseline severity.

\textbf{Inputs.} The E1 and E2 predictions.

\textbf{Outputs.}
\begin{itemize}
\item \texttt{runs/experiments/E9/adversarial\_robustness.json}.
\item \texttt{runs/experiments/E9/manifest.json}.
\end{itemize}

\textbf{Feeds.} Table 38. The reported $96.5\%$ robustness is the
ratio of adaptive-blend accuracy to baseline accuracy, not the raw
accuracy itself.

\subsection{E10 --- Statistical significance}

\textbf{Purpose.} Run the four significance tests (paired bootstrap,
Wilcoxon signed-rank, permutation, Friedman) over the E1--E9
artefacts.

\textbf{Inputs.} The E1, E2, E3, E5, E6, E7, E8 and E9 manifests and
prediction files.

\textbf{Outputs.}
\begin{itemize}
\item \texttt{runs/experiments/E10/statistical\_analysis.json}.
\item \texttt{runs/experiments/E10/significance\_table.tex}.
\end{itemize}

\textbf{Feeds.} Table 39. E10 depends on every other experiment, so
it is the last node in the orchestrator's DAG.

\newpage

% latex/supplement/figures.tex
%
% Figure captions for the seven data figures the extended suite
% reproduces. Each caption records what the figure shows, what the
% underlying data is, and the specific finding it supports. The
% captions are the reference; the figure itself is generated by
% extended_experimentation/reporting/a3e_figures.py.

\section{Figure captions}

The seven figures reproduced by
\texttt{extended\_experimentation/reporting/a3e\_figures.py} are
described below. A reviewer regenerating them from the pipeline
should see the same qualitative shapes the thesis reports; small
numerical differences in the plotted points are expected and are
covered by the tolerance bands in the evaluation suite.

\subsection*{Figure F1 --- Detection accuracy vs.\ attack severity}

\emph{Caption.} Detection accuracy across the ten-scenario attack
catalogue, plotted against severity. Each point is one scenario at one
severity variant; the horizontal dashed line marks the
$99.94\%$ thesis reference. Accuracy remains above $99.4\%$ for all
scenarios up to severity $0.65$; the adaptive-blend variant of the
tunnel scenario (S2) at severity $0.65$ is the first to fall below the
reference, at $99.31\%$.

\emph{Supports.} Section~6.2.2, Table~30.

\subsection*{Figure F2 --- Byzantine robustness vs.\ Byzantine fraction}

\emph{Caption.} Mean and maximum attack-to-reference ratios across
the Byzantine sweep. The mean ratio (blue) rises from $0.42$ at
fraction $0.05$ to $1.87$ at fraction $0.40$; the maximum ratio
(orange) rises from $0.58$ to $2.94$. The gap between the two curves
is the heterogeneity across the eight attack strategies; the mean is
what the aggregate-accuracy panel of Table~31 reports, and the
maximum is what the worst-case robustness discussion in
Section~6.3.4 cites.

\emph{Supports.} Section~6.3, Table~31.

\subsection*{Figure F3 --- DP epsilon vs.\ utility}

\emph{Caption.} Detection accuracy (green) and false-positive rate
(purple) as a function of the differential-privacy budget $\eps$. The
vertical dashed line marks the thesis's $\eps = 0.68$. At
$\eps = 0.68$ the accuracy loss relative to the non-private baseline
is $0.4\%$ and the FPR loss is $0.05\%$; both are within the
tolerances recorded in the evaluation suite. The sharp utility drop
below $\eps = 0.4$ is the reason the thesis chooses $\eps = 0.68$
rather than a smaller budget.

\emph{Supports.} Section~6.6.2, Table~32.

\subsection*{Figure F4 --- Latency CDF by protocol}

\emph{Caption.} Latency percentiles across DoH, DoT and DoQ. The
curves plot the $P_{50}$, $P_{95}$ and $P_{99}$ points of each
protocol's latency distribution. DoQ is the lowest-latency transport
at $P_{50}$ ($0.81\,\mathrm{ms}$) because its UDP-based transport
avoids the TCP handshake; DoH is the highest ($1.12\,\mathrm{ms}$) for
the same reason. The horizontal dashed line marks the thesis's
$1.42\,\mathrm{ms}$ P99 reference, which the pooled distribution
meets.

\emph{Supports.} Section~6.7.1, Table~35.

\subsection*{Figure F5 --- Throughput vs.\ concurrency}

\emph{Caption.} Peak QPS as a function of concurrency. Throughput
rises linearly up to concurrency $= 32$ and flattens thereafter as the
pipeline's inference path saturates. The horizontal dashed line
marks the $22.5\,\mathrm{M}\,\mathrm{QPS}$ thesis reference, which
corresponds to the concurrency $= 32$ point.

\emph{Supports.} Section~6.7.2, Table~35.

\subsection*{Figure F6 --- QHE NTT vs.\ convolution speedup}

\emph{Caption.} Median transform time for the negacyclic NTT
(blue) versus naive polynomial multiplication (orange). The
annotated speedup factor is the ratio of the two medians; on the
reference machine the NTT is $18.4\times$ faster. The naive
implementation is the C reference from
\texttt{c/reference/src/ntt.c}; the NTT is the Rust implementation
from \texttt{rust/a3e\_core/src/ntt.rs}.

\emph{Supports.} Section~6.5.1, Table~33.

\subsection*{Figure F7 --- Multi-protocol accuracy}

\emph{Caption.} Detection accuracy for DoH, DoT and DoQ under the
same detector. All three protocols exceed the $99.94\%$ thesis
reference; the difference between them is below the tolerance band,
which is the evidence for the claim that the detector generalises
across transports.

\emph{Supports.} Section~6.7.3, Table~37.

\newpage

% ----- Generated result tables ------------------------------------- %
\section{Result tables}

The tables below are generated by
\texttt{extended\_experimentation/reporting/a3e\_result\_tables.py}
from the experiment manifests under
\texttt{runs/experiments/}. They are \verb|\input{}| here rather than
inlined so that a change to the generators updates this section
without a manual edit.

\subsection*{Table S1 --- Per-scenario breakdown}
\input{tables/table_S1}

\subsection*{Table S2 --- Byzantine robustness sweep}
\input{tables/table_S2}

\subsection*{Table S3 --- Differential-privacy epsilon sweep}
\input{tables/table_S3}

\subsection*{Table S4 --- QHE primitive costs}
\input{tables/table_S4}

\subsection*{Table S5 --- Selective-operation accuracy}
\input{tables/table_S5}

\subsection*{Table S6 --- Latency percentiles}
\input{tables/table_S6}

\subsection*{Table S7 --- Ablation summary}
\input{tables/table_S7}

\subsection*{Table S8 --- Multi-protocol accuracy}
\input{tables/table_S8}

\subsection*{Table S9 --- Adversarial robustness}
\input{tables/table_S9}

\newpage

% ----- Appendix: how to reproduce ---------------------------------- %
\appendix
\section{Reproducing the supplement}

The supplement is regenerated by running the pipeline end-to-end:

\begin{verbatim}
    make build
    bash scripts/run_pipeline.sh runs/$(date +%Y%m%d)
    python3 -m extended_experimentation.cli tables \
        --manifests-root runs/*/experiments \
        --out latex/supplement/tables
    pdflatex -output-directory latex/supplement/build \
        latex/supplement/main.tex
\end{verbatim}

\noindent
The first command builds every component (Python, Rust, Go, C); the
second runs the ten experiments; the third regenerates the S1--S9
tables; the fourth typesets this supplement. If any step fails, the
subsequent step refuses to proceed, so a successful \texttt{pdflatex}
is evidence that the whole pipeline ran.

\section{What the supplement does not contain}

The supplement deliberately omits:

\begin{itemize}
\item The schematic architecture figures from the thesis body. Those
      are drawn in TikZ and belong in the thesis, not in a data
      supplement.
\item The proof sketches for the cryptographic soundness claims.
      Those are in the thesis body, Chapter~6, Section~6.5.3.
\item The full per-window prediction data. That is available as
      JSONL under \texttt{runs/experiments/*/predictions.jsonl} and as
      SQL rows in \texttt{runs/results.db}; inlining it here would
      have produced a 500-page document whose value is the ability
      to grep rather than to read.
\end{itemize}

\noindent
The supplement is the human-readable projection of the machine-
readable artefacts. A reviewer who needs to see a prediction has the
artefact; a reviewer who needs to see a summary has the supplement.

\end{document}
